\documentclass[10pt,a4paper]{amsart}
\usepackage[margin=19mm]{geometry}
\usepackage{amsmath,amssymb,mathtools}
\usepackage[scr=rsfs,cal=euler]{mathalfa}
\usepackage{xfrac}
\usepackage[unicode,hidelinks,bookmarks=false]{hyperref}
\newcommand{\ie}{i.e.\ }
\newcommand{\cf}{cf.\ }
\newcommand{\resp}{respectively\ }
\newcommand{\Subsection}{\subsection}
\providecommand{\nobreakdash}{\nobreak\hbox{-}}
\setlength{\emergencystretch}{2em}
\multlinegap0pt
\usepackage{amstext}
\usepackage{mathrsfs}
\usepackage{bm}
\usepackage{mathtools}
\usepackage{color}
\usepackage{mathbbol}
\usepackage{multirow}
\usepackage{enumitem}
\usepackage{tabularx}
\usepackage{blkarray}
\usepackage{placeins}
\usepackage{array}
\usepackage{booktabs}
\usepackage{xy}
\usepackage{tikz}
\usepackage{tikz-cd}
\newcommand{\xycenter}[1]{
	\begin{center}
	\mbox{\xymatrix{#1}}
	\end{center}
	}
\newtheorem{theorem}{Theorem}[section]
\newtheorem*{expectation}{Expectation}
\newtheorem{proposition}[theorem]{Proposition}
\newtheorem{lemma}[theorem]{Lemma}
\newtheorem{corollary}[theorem]{Corollary}
\newtheorem{problem}[theorem]{Problem}
\numberwithin{equation}{section}
\newtheorem{theoremintro}{Theorem}
\newtheorem{definition}[theorem]{Definition}
\newtheorem{notation}[theorem]{Notation}
\newtheorem{convention}[theorem]{Convention}
\newtheorem{example}[theorem]{Example}
\newtheorem{conjecture}[theorem]{Conjecture}
\newtheorem{construction}[theorem]{Construction}
\newtheorem*{maintheorem}{Theorem \ref{tSaturatedPrelogY}}
\newtheorem{claim}[equation]{Claim}
\newtheorem{remark}[theorem]{Remark}
\newcommand{\io}[2]{\iota_{\{#1\}>\{#2\}}}
\newcommand{\ioUpper}[2]{\iota^*_{\{#1\}>\{#2\}}}
\newcommand{\ioLower}[2]{\iota_{\{#1\}>\{#2\}*}}
\newcommand{\vin}{\rotatebox[origin=c]{-90}{$\in$}}
\newcommand{\sheaf}[1]{\mathscr{#1}}
\newcommand{\DD}{\sheaf{D}}
\newcommand{\LL}{\sheaf{L}}
\newcommand{\OO}{\sheaf{O}}
\newcommand{\QQ}{\sheaf{Q}}
\newcommand{\MM}{\sheaf{M}}
\newcommand{\EE}{\sheaf{E}}
\newcommand{\FF}{\sheaf{F}}
\newcommand{\II}{\sheaf{I}}
\newcommand{\HH}{\sheaf{H}}
\newcommand{\NN}{\sheaf{N}}
\newcommand{\VV}{\sheaf{V}}
\newcommand{\ZZ}{\sheaf{Z}}
\newcommand{\WW}{\sheaf{W}}
\newcommand{\BB}{\sheaf{B}}
\newcommand{\YY}{\sheaf{Y}}
\newcommand{\XX}{\sheaf{X}}
\newcommand{\sA}{\sheaf{A}}
\newcommand{\sS}{\sheaf{S}}
\newcommand{\UU}{\sheaf{U}}
\newcommand{\perfect}[1]{{\color{black} #1}}
\newcommand{\formalbase}{\mathrm{Spec}\,\mathbb{C}[[t]]}
\newcommand{\valuations}[1]{\mathsf{#1}}
\newcommand{\residue}{\partial}
\newcommand{\divisor}{\mathrm{div}}
\newcommand{\Brtwo}{{}_2\mathrm{Br}}
\newcommand{\Pictwo}{{}_2\mathrm{Pic}}
\newcommand{\isom}{\cong}
\newcommand{\isomto}{\simto}
\newcommand{\isometry}{\cong}
\newcommand{\Z}{\mathbb Z}
\newcommand{\N}{\mathbb N}
\newcommand{\A}{\mathbb A}
\newcommand{\C}{\mathbb C}
\renewcommand{\P}{\mathbb P}
\newcommand{\Q}{\mathbb Q}
\newcommand{\Gm}{\mathbb{G}_{\mathrm{m}}}
\newcommand{\Gms}[1]{\mathbb{G}_{\mathrm{m},#1}}
\DeclareMathOperator{\AAut}{\mathbf{Aut}}
\DeclareMathOperator{\Aut}{\mathrm{Aut}}
\DeclareMathOperator{\Br}{\mathrm{Br}}
\DeclareMathOperator{\Isom}{\mathrm{Isom}}
\DeclareMathOperator{\IIsom}{\Group{Isom}}
\DeclareMathOperator{\PPic}{\sheaf{P}\!\mathit{ic}}
\DeclareMathOperator{\Pic}{\mathrm{Pic}}
\DeclareMathOperator{\rk}{\mathrm{rk}}
\DeclareMathOperator{\SSpec}{\mathbf{Spec}}
\DeclareMathOperator{\Spec}{\mathrm{Spec}}
\DeclareMathOperator{\EExt}{\sheaf{E}\!\mathit{xt}}
\DeclareMathOperator{\Ext}{\mathrm{Ext}}
\DeclareMathOperator{\coker}{\mathrm{coker}}
\DeclareMathOperator{\Chow}{\mathrm{CH}}
\DeclareMathOperator{\Num}{\mathrm{Num}}
\DeclareMathOperator{\Chownum}{\mathrm{Num}}
\DeclareMathOperator{\Chowprelog}{\mathrm{CH}^*_{\mathrm{prelog}}}
\DeclareMathOperator{\Chowprelognum}{\mathrm{Num}^*_{\mathrm{prelog}}}
\DeclareMathOperator{\Chowprelogsat}{\mathrm{Chow}_{\mathrm{prelog, sat}}}
\DeclareMathOperator{\Numsatprelog}{\mathrm{Num}_{\mathrm{prelog, sat}}}
\DeclareMathOperator{\Hilb}{Hilb}
\DeclareMathOperator{\codim}{codim}
\DeclareMathOperator{\SL}{SL}
\newcommand{\inv}{^{-1}}
\newcommand{\sep}{^{\mathrm{s}}}
\newcommand{\mult}{^{\times}}
\newcommand{\dual}{^{\vee}}
\newcommand{\tensor}{\otimes}
\newcommand{\bslash}{\smallsetminus}
\newcommand{\mapto}[1]{\xrightarrow{#1}}
\newcommand{\ol}[1]{\overline{#1}}
\newcommand{\ul}[1]{\underline{#1}}
\newcommand{\wt}[1]{\widetilde{#1}}
\newcommand{\et}{\mathrm{\acute{e}t}}
\newcommand{\linedef}[1]{\textsl{#1}}
\newcommand{\Het}{H_{\et}}
\newcommand{\ur}{\mathrm{nr}}
\newcommand{\Hur}{H_{\ur}}
\newcommand{\Brur}{\Br_{\ur}}
\newcommand{\merk}{\mathrm{r}}
\newcommand{\Hr}{H_{\merk}}
\newcommand{\Pfister}[1]{\ll\!{#1}\gg}
\newcommand{\quadform}[1]{<\! #1 \!>}
\newcommand{\Local}{\mathsf{Local}}
\newcommand{\Ab}{\mathsf{Ab}}
\newcommand{\Var}{\mathsf{Var}}
\newcommand{\Frac}{\mathrm{Frac}}
\newcommand{\im}{\mathrm{im}}
\newcommand{\CH}{\mathrm{CH}}
\newcommand{\CM}{\mathsf{CM}}
\newcommand{\res}{\mathrm{res}}
\newcommand{\cores}{\mathrm{cor}}
\newcommand{\ord}{\mathrm{ord}}
\newcommand{\Norm}{\mathrm{N}}
\newcommand{\vp}{\pi}
\newcommand{\Hom}{\mathrm{Hom}}
\newcommand{\id}{\mathrm{id}}
\newcommand{\diag}{\mathrm{diag}}
\newcommand{\Pfad}{\mathrm{adj}^{\mathrm{Pf}}}
\newcommand{\Pf}{\mathrm{Pf}}
\newcommand{\adj}{\mathrm{adj}}
\newcommand{\wMM}{\widetilde{\MM}}
\newcommand{\bMM}{\overline{\mathcal{M}}}
\newcommand{\oMM}{\mathcal{M}}
\newcommand{\mf}{\mathcal{MF}}
\newcommand{\mfss}{\mf^{ss}}
\newcommand{\mfps}{\mf^{ps}}
\newcommand{\sk}{\mathcal{S}}
\newcommand{\skss}{\sk^{ss}}
\newcommand{\skps}{\sk^{ps}}
\newcommand{\wsk}{\widetilde{\sk}}
\newcommand{\wskss}{\widetilde{\sk}^{ss}}
\newcommand{\wskps}{\widetilde{\sk}^{ps}}
\newcommand{\GL}{\mathrm{GL}}
\newcommand{\GLsix}{\GL_6 (\C )}
\newcommand{\LLF}{\overline{\mathcal{B}}}
\newcommand{\wLLF}{\widetilde{\mathcal{B}}}
\newcommand{\CF}{\overline{\mathcal{A}}}
\newcommand{\wCF}{\widetilde{\mathcal{A}}}
\newcommand{\G}{\mathbb{G}}
\newcommand{\cubicsexplicit}{\P\bigl(H^0\bigl(\OO_{\P^4}(3)\bigr)\bigr)}
\newcommand{\cubics}{\mathcal{C}}
\newcommand{\Conics}{\Hilb_{\mathrm{conics}}}
\newcommand{\Skewlines}{\Hilb_{\mathrm{skewlines}}}
\newcommand{\ring}{R}
\newcommand{\red}[1]{{\color{red} #1}}
\newcommand\WHY{{\color{red}\textsf{WHY?}}~}
\renewcommand\theenumi{\it\alph{enumi}}
\renewcommand\labelenumi{\theenumi)}
\newcommand*{\DashedArrow}[1][]{\mathbin{\tikz [baseline=-0.25ex,-latex, dashed,#1] \draw [#1] (0pt,0.5ex) -- (2.3em,0.5ex);}}
\newcommand{\bP}{\mathbb{P}}
\newcommand{\bA}{\mathbb{A}}
\newcommand{\bH}{\mathbb{H}}
\newcommand{\bR}{\mathbb{R}}
\newcommand{\bG}{\mathbb{G}}
\newcommand{\bQ}{\mathbb{Q}}
\newcommand{\bZ}{\mathbb{Z}}
\newcommand{\bF}{\mathbb{F}}
\newcommand{\bC}{\mathbb{C}}
\newcommand{\bM}{\mathbb{M}}
\newcommand{\bL}{\mathbb{L}}
\newcommand{\bB}{\mathbb{B}}
\newcommand{\fX}{\mathfrak{X}}
\newcommand{\sH}{\mathrm{H}}
\newcommand{\sT}{\mathrm{T}}
\newcommand{\calA}{\mathcal{A}}
\newcommand{\calC}{\mathcal{C}}
\newcommand{\calE}{\mathcal{E}}
\newcommand{\calF}{\mathcal{F}}
\newcommand{\calK}{\mathcal{K}}
\newcommand{\calH}{\mathcal{H}}
\newcommand{\calN}{\mathcal{N}}
\newcommand{\calT}{\mathcal{T}}
\newcommand{\calh}{\mathfrak{h}}
\newcommand{\fN}{\mathfrak{N}}
\newcommand{\calM}{\mathcal{M}}
\newcommand{\calO}{\mathcal{O}}
\newcommand{\calL}{\mathcal{L}}
\newcommand{\calP}{\mathcal{P}}
\newcommand{\calI}{\mathcal{I}}
\newcommand{\calX}{\mathcal{X}}
\newcommand{\calZ}{\mathcal{Z}}
\newcommand{\calB}{\mathcal{B}}
\newcommand{\calY}{\mathcal{Y}}
\newcommand{\calD}{\mathcal{D}}
\newcommand{\calJ}{\mathcal{J}}
\newcommand{\Der}{\mathrm{Der}}
\newcommand{\Bir}{\mathrm{Bir}}
\newcommand{\Stab}{\mathrm{Stab}}
\newcommand{\Span}{\mathrm{span}}
\newcommand{\Sym}{\mathrm{Sym}}
\newcommand{\PGL}{\mathrm{PGL}}
\newcommand{\Mov}{\mathrm{Mov}}
\newcommand{\Sat}{\mathrm{Sat}}
\newcommand{\Sp}{\mathrm{Sp}}
\newcommand{\Proj}{\mathrm{Proj}}
\newcommand{\Sing}{\mathrm{Sing}}
\newcommand{\s}{\mathrm{sp}}
\newcommand{\Ker}{\mathrm{Ker}}
\newcommand{\Gr}{\mathrm{Gr}}
\newcommand{\Sec}{\mathrm{Sec}}
\newcommand{\corank}{\mathrm{corank}}
\newcommand{\rank}{\mathrm{rank}}
\newcommand{\rad}{\mathrm{rad}}
\newcommand{\Nef}{\mathrm{Nef}}
\newcommand{\Lin}{\mathrm{Lin}}
\newcommand{\NS}{\mathrm{NS}}
\newcommand{\Def}{\mathrm{Def}}
\newcommand{\SO}{\mathrm{SO}}
\newcommand{\Bis}{\mathrm{Bis}}
\newcommand{\Prym}{\mathrm{Prym}}
\newcommand{\Mon}{\mathrm{Mon}}
\newcommand{\gquot}{/\!\!/}
\newcommand{\spec}{\mathrm{Spec}\;} 
\newcommand{\sslash}{\mathbin{/\mkern-6mu/}}
\newcommand{\git}{/\kern-0.2em/}
\newcommand{\Lneg}{L_{-}}
\newcommand{\Linv}{L_+}
\newcommand{\pforms}{\Omega^{[p]}}
\newcommand{\OG}{\Lambda}
\newcommand{\markman}{\mathcal{W}^{pex}}
\newcommand{\mongardi}{\mathcal{W}}
\newcommand{\leech}{\mathbb{L}}
\newcommand{\borcherds}{\mathbb{B}}
\newcommand{\conway}{\textnormal{Co}_0}
\newcommand{\coinv}{\Lambda_G}
\newcommand{\Ghash}{G^{\#}}
\newcommand{\sat}{\textnormal{Sat}}
\newcommand{\cubic}{\mathbb{C}}
\newcommand{\tH}{\widetilde{\mathrm{H}}}
\newcommand{\Db}{\mathrm{D}^{\mathrm{b}}}
\newcommand{\Ku}{\mathrm{Ku}}
\setcounter{MaxMatrixCols}{24}
\DeclareMathOperator{\divi}{div}
\newcommand{\stevell}[1]{{\color{brown} \sf $\clubsuit\clubsuit\clubsuit$ Stevell: [#1]}}
\newcommand\sfedit[1]{{\color{brown} \sf #1}}
\newcommand{\Stevell}[1]{\normalsize{{\color{red}\footnote{{\color{blue}#1}}}{\marginpar[{\color{red}\hfill\tiny\thefootnote$\rightarrow$}]{{\color{red}$\leftarrow$\tiny\thefootnote}}}}}
\newcommand{\christian}[1]{{\color{blue} \sf $\clubsuit\clubsuit\clubsuit$ Christian: [#1]}}
\newcommand\cbedit[1]{{\color{blue} \sf #1}}
\newcommand{\Christian}[1]{\normalsize{{\color{red}\footnote{{\color{blue}#1}}}{\marginpar[{\color{red}\hfill\tiny\thefootnote$\rightarrow$}]{{\color{red}$\leftarrow$\tiny\thefootnote}}}}}
\newcommand{\lisa}[1]{{\color{purple} \sf $\clubsuit\clubsuit\clubsuit$ Lisa: [#1]}}
\newcommand\lmedit[1]{{\color{purple} \sf #1}}
\newcommand{\Lisa}[1]{\normalsize{{\color{red}\footnote{{\color{purple}#1}}}{\marginpar[{\color{red}\hfill\tiny\thefootnote$\rightarrow$}]{{\color{red}$\leftarrow$\tiny\thefootnote}}}}}
\begin{document}
\title[Naive atoms of blowups: examples]{Naive atoms of blowups: examples}
\author{Christian Böhning, Hans-Christian Graf von Bothmer, Zac Su'a}
\date{}
\maketitle
\noindent\emph{Source and licence.} Naive atoms of blowups: examples. Version arXiv:2606.17884v2 (CC BY 4.0). This material is distributed under the Creative Commons Attribution 4.0 International licence, http://creativecommons.org/licenses/by/4.0/, as recorded in the arXiv OAI licence metadata for this exact version and shown on the abstract page https://arxiv.org/abs/2606.17884v2. Full rights notice: licenses/arxiv-2606.17884-CC-BY-4.0.txt.\par\smallskip
\noindent\textit{Editorial scope.} Counted unit: the original \texttt{theorem} environment \texttt{tMatrixCubicFourfoldLine} at source lines 1127--1185 together with its complete proof at lines 1187--1207; the delivered document keeps source lines 1068--1208 so that every referenced label and every citation resolves inside it. Native editable source is the paper's own concatenated TeX; the AMS wrapper is editorial formatting only. No publisher class, upstream script or unrelated artwork is redistributed.
\section{Cubic fourfolds blown up in a line}\label{sLine}

Let $X$ be a cubic fourfold, and let $\widetilde{X}$ be the blow up of $X$
along a line $L \subset X$. We denote by $E$ the exceptional divisor. Let
$P$ be the class of a fibre of the projection $E \to L$, and let $F$ be the
class of a line in such a fibre. Finally, we set
\[
S := -E^2,
\]
which is the class of a relative line.

\begin{proposition}\label{pIntersectionRing}
The intersection ring of $\widetilde{X}$ is generated, as a $\Z$-module,
by the classes
\[
    1,H,H^2,L,pt,E,P,S,F
\]
and its product structure is determined by the relations
\begin{gather*}
    H^3 = 3L, \qquad
    \mathrm{pt} = H L, \qquad
    P = E H, \qquad
    S = -E^2, \qquad
    F = S H, \qquad
    H^2E = 0, \\
    (H-E)^3 = 2(L-F).
\end{gather*}
\end{proposition}


\begin{proof}
The first five relations are definitions. Moreover, $H^2E=0$, since $H^2$
is represented by a codimension-two linear section which does not meet the
blowup center $L$.

It remains to prove the last relation. Consider the projection from the line
$L$ to $\P^3$. After blowing up $L$, this projection is given by the
linear system $|H-E|$ on $\widetilde{X}$. Hence the class of a fibre is $(H-E)^3$.

On the other hand, such a fibre is the transform of the residual curve to
$L$ in the intersection of $X$ with a plane containing $L$. This residual
curve is a conic, since $X$ is cubic, and it meets $L$ in two points.
Therefore its strict transform has class $2L - 2F$. The last relation follows. 
\end{proof}

Notice that the monoid of effective curve classes on $\widetilde{X}$ is generated by $F$ and $L-F$. 

\begin{lemma}\label{lDualBasis}
The dual basis to the basis  $\{1,H,H^2,L,pt,E,P,S,F\}$ is given by
\[
\{ pt,L,\frac{1}{3}H^2,H,1,-F,-S-P,-P,-E\}. 
\]
\end{lemma}

\begin{proof}
This is a direct computation.
\end{proof}



\begin{theorem}\label{tMatrixCubicFourfoldLine}
Let $\widetilde{X}$ be the blowup of $X$ along a line $L$. As before, set $q = q^L$ and $t = (q^F)^{-1}$.
Then the matrix $M_{(\widetilde{X},-K_{\widetilde{X}})}$ representing quantum
multiplication by $-K_{\widetilde{X}}$ is given by
\[
\renewcommand{\arraystretch}{1.2}
\begin{array}{c|ccccc|cccc}
	& pt & L & \frac{1}{3}H^2 & H & \widetilde{X} & -F & -S-P & -P & -E \\
	\hline
       \widetilde{X} & 0&3&0&0&0&-2&0&0&0\\
       H & 4q^{2}t^{2}&5qt&3&0&0&-5qt&-2&0&0\\
       H^2 & 18q&6q^{2}t^{2}&7qt&9&0&-6q^{2}t^{2}&-15qt&-6qt&0\\
       L & 56q^{2}t&15q&2q^{2}t^{2}&5qt&3&0&-4q^{2}t^{2}&-2q^{2}t^{2}&-5qt\\
       pt & 80q^{3}t^{2}&56q^{2}t&6q&4q^{2}t^{2}&0&-24q^{2}t&0&0&-4q^{2}t^{2}\\
       \hline
       E & 4q^{2}t^{2}&5qt&0&0&0&-5qt&3&2&0\\
       P & 0&2q^{2}t^{2}&2qt&0&0&-2q^{2}t^{2}&-5qt&-qt&2\\
       S & 0&2q^{2}t^{2}&3qt&2&0&\frac{-2q^{2}t^{3}-2}{t}&-5qt&-4qt&1\\
       F & 24q^{2}t&0&2q^{2}t^{2}&5qt&2&3q&\frac{-4q^{2}t^{3}-2}{t}&-2q^{2}t^{2}&-5qt
 \end{array}
\]
The characteristic polynomial is
\[
    \chi =
    \frac{(2qt+\lambda)^2 F_7}{t^2},
\]
where $F_7$ is an irreducible polynomial of degree $7$. 

The coefficient of the leading term of $\chi$ with respect to $t$ is
\[
    -16\lambda^2(\lambda^3 - 729q),
\]
which recovers the characteristic polynomial of the cubic fourfold $X$.

After the change of variables
\[
    \lambda \mapsto \lambda + \frac{2i}{\sqrt{t}},
    \qquad
    q \mapsto \frac{q}{i\sqrt{t}},
\]
the coefficient of the leading term with respect to $\sqrt{t}$ is
\[
    512(\lambda^2 - 4q),
\]
which recovers the characteristic polynomial of the blow-up center $L$.

Similarly, after the change of variables
\[
    \lambda \mapsto \lambda - \frac{2i}{\sqrt{t}},
    \qquad
    q \mapsto \frac{q}{i\sqrt{t}},
\]
the coefficient of the leading term with respect to $\sqrt{t}$ is
\[
    -512(\lambda^2 + 4q),
\]
which recovers the characteristic polynomial of the blow-up center $L$, up to
the sign change $q \mapsto -q$.
\end{theorem}

\begin{proof}
Apart from the commutativity relation $H \star E = E \star H$, the following
Gromov--Witten invariants were used as input for the computation of
$M_{(\widetilde{X},-K_{\widetilde{X}})}$ in Macaulay2:
\begin{align*}
    \langle L \rangle_{L-F}
        &= 5
        && \text{(lines on a cubic surface meeting two given lines)}, \\
    \langle \mathrm{pt} \rangle_{2L-2F}
        &= 1
        && \text{(planes through $L$ containing a general point)}, \\
    \langle H^2,\mathrm{pt}\rangle_L
        &= 6
        && \text{(degree of the cone swept out by the lines through a general point)}.
\end{align*}

The characteristic polynomial, the changes of variables, and the leading
terms in the corresponding limits were also computed using Macaulay2 \cite{BvBS26}.

Finally, the absolute irreducibility of the polynomial $F_7$ was checked using SageMath, {\ttfamily https://sagecell.sagemath.org}.
\end{proof}


\begin{thebibliography}{99}
\bibitem[BvBS26]{BvBS26} B\"ohning, Chr., Graf von Bothmer, H.-C.., Su'a, Z., \emph{Macaulay 2 files for ``Naive atoms of blowups: examples"}, available at \href{https://zenodo.org/uploads/20625923}{https://zenodo.org/uploads/20625923}
\end{thebibliography}
\end{document}
